\documentclass[12pt,letterpaper]{article}
\usepackage[utf8]{inputenc}
\usepackage[T1]{fontenc}
\usepackage[spanish]{babel}
\usepackage{amsmath, amssymb}
\usepackage{geometry}
\usepackage{enumitem}
\usepackage{graphicx}

\geometry{margin=2.5cm}

\title{\textbf{Ejercicios tipo -- Certamen 3 -- Microeconomía I ICS-162}}
\date{}

\begin{document}

\maketitle

%=============================================================


%------------------------------------------------------------
\subsection*{Pregunta 1}

Amanda afirma que su situación como emprendedor es muy desventajosa debido a que debe arrendar un local para su carnicería. Mientras que su competidora, Agustina, es propietaria del local donde tiene instalada su carnicería. Según Amanda, evidentemente, Agustina tiene mayores beneficios económicos que ella. Comente esta afirmación.


%------------------------------------------------------------
\subsection*{Pregunta 2}

La firma de tornillos Fantásticos S.A. está en una industria en competencia perfecta. Una persona que no ha tomado microeconomía escucha que tiene beneficios normales y pregunta: ``¿qué es eso?'' Su compañero que sí tomó microeconomía le dice que quiere decir que los beneficios son iguales a cero. Entonces exclama: ``¡Entonces para qué existe si no gana nada! Eso no es normal''. ¿Qué opina de esta afirmación?



%------------------------------------------------------------
\subsection*{Pregunta 3}

Federico asesora una empresa que debe decidir si cerrar o no. Él ve que tiene beneficios subnormales y lo que gana no alcanza a cubrir sus costos fijos. ¿Es esta suficiente información para tomar la decisión? Si tu respuesta es positiva indica qué conviene en este caso y por qué; si es negativa, qué información faltaría.



%------------------------------------------------------------
\subsection*{Pregunta 4}

Una empresa agrícola opera en un mercado perfectamente competitivo de producción de trigo. En el corto plazo, obtiene beneficios económicos positivos. Explique, utilizando conceptos vistos en clases, qué ocurrirá con el precio del trigo, la cantidad producida por la empresa individual y el beneficio económico en el largo plazo. ¿Cuál será la situación de equilibrio de largo plazo para esta empresa?

%------------------------------------------------------------
\subsection*{Pregunta 5}

Analice la siguiente afirmación: \textit{``Si la elasticidad del costo con respecto a la producción es menor que 1, entonces existen economías de escala.''} Justifique su respuesta y apoye su explicación con un gráfico.



%----------------------------------------------------------


%=============================================================
\newpage


%------------------------------------------------------------
\subsection*{Ejercicio 1}

Una empresa en competencia perfecta en el corto plazo tiene una curva de costo total:
\[
CT = 500 + q^3 - 20q^2 + 110q
\]

\subsubsection*{(a) (10 pts) Calcule el precio de cierre de la firma en el corto plazo.}


\subsubsection*{(b) (5 pts) Calcule los beneficios en ese punto e indique si son normales, subnormales o sobrenormales.}

\subsubsection*{(c) (5 pts) Formule la curva de oferta (u oferta inversa) de la firma.}


%------------------------------------------------------------
\subsection*{Ejercicio 2}

El maíz es producido en condiciones de competencia perfecta. Cada productor individual tiene una curva de costos medios a largo plazo con forma de U, alcanzando su mínimo en \$4 por tonelada cuando produce 500 toneladas.

\subsubsection*{(a) (12 pts) Determine el precio de equilibrio de largo plazo y las toneladas demandadas en total.}



\subsubsection*{(b) (8 pts) ¿Cuántos productores habrá en el mercado y cuál será la utilidad de cada uno?}

\end{document}